\section{Method}

The staged method builds a deterministic spectral-style coordinate, sorts units
by that coordinate, and evaluates population-balanced sweep cuts. Recursive calls pass
proportional target fractions so odd and non-power-of-two district counts are
handled without silently changing the intended seat allocation.

In the atlas framing, the candidates are not arbitrary graph cuts. They are
prefix cuts induced by the Fiedler ordering. The ordering explains why the
candidate set exists; the sweep table explains why a particular prefix is
selected or skipped. The downstream consequence of the selected root cut is a
pair of recursive workloads with explicit left/right seat counts.

\subsection{Reproducible Procedure}

The constructor operates on the tract adjacency graph supplied by the BISECT
state pipeline. Vertices carry population weights, edges encode adjacency, and
the current implementation treats edge weights as the graph weights available
from that pipeline. For each recursive subproblem:

\begin{enumerate}
  \item Build the weighted graph induced by the current unit set.
  \item Build a centered anchor vector from stable unit order and smooth it
        over adjacency for the configured iteration budget.
  \item Normalize the coordinate after each smoothing pass and sort units by
        coordinate, breaking ties by stable unit identifier.
  \item Sweep prefixes in that order, reject cuts outside the declared
        population tolerance, and select the feasible cut with minimum edge cut,
        breaking remaining ties by sweep index.
  \item Recurse on the two induced subgraphs with the requested left and right
        seat counts until each leaf represents one district.
\end{enumerate}

This procedure makes determinism an implementation contract: the same graph,
weights, target fractions, and tie-breaking identifiers must produce the same
assignment and lineage. The graph assumptions are deliberately narrow. Water
gaps, islands, disconnected components, and alternative adjacency definitions
belong to the input graph construction layer, and near-symmetric graphs can
make the spectral coordinate weakly informative even when the run remains
deterministic.

\subsection{Odd-Seat Split Sketch}

For non-power-of-two district counts, the important visual is the target tree,
not a ratio label. A seven-district run first asks the root cut to create
workloads for three and four seats. The four-seat side can then split evenly,
while the three-seat side carries one odd split. The constructor therefore
chooses each local bisection with the population target implied by the child
seat counts, not by a universal 50/50 rule.

\begin{table}[h]
\centering
\small
\begin{tabular}{p{0.12\linewidth}p{0.18\linewidth}p{0.32\linewidth}p{0.27\linewidth}}
\toprule
Depth & Node target & Requested child targets & Local target fraction \\
\midrule
0 & 7 seats & 3 seats / 4 seats & \(3/7\) then \(4/7\) \\
1 & 3 seats & 1 seat / 2 seats & \(1/3\) then \(2/3\) \\
1 & 4 seats & 2 seats / 2 seats & \(1/2\) then \(1/2\) \\
2 & 2 seats & 1 seat / 1 seat & \(1/2\) then \(1/2\) \\
\bottomrule
\end{tabular}
\caption{Recursive target propagation for \(k=7\). Each row describes the
population target used by a local spectral sweep, so odd-seat recursion remains
auditable rather than implicit.}
\end{table}
